\documentclass[11pt,reqno]{amsart}

\usepackage{amsmath,amssymb,amsthm}
\usepackage[margin=1.1in]{geometry}
\usepackage{microtype}
\usepackage[colorlinks=true,linkcolor=blue,citecolor=blue]{hyperref}

\theoremstyle{plain}
\newtheorem{theorem}{Theorem}[section]
\newtheorem{proposition}[theorem]{Proposition}
\newtheorem{lemma}[theorem]{Lemma}
\newtheorem{corollary}[theorem]{Corollary}
\newtheorem{conjecture}[theorem]{Conjecture}
\theoremstyle{definition}
\newtheorem{definition}[theorem]{Definition}
\newtheorem{remark}[theorem]{Remark}
\newtheorem{example}[theorem]{Example}

\newcommand{\R}{\mathbb{R}}
\newcommand{\Z}{\mathbb{Z}}
\newcommand{\Q}{\mathbb{Q}}
\newcommand{\C}{\mathbb{C}}
\newcommand{\SU}{\mathrm{SU}(2)}
\newcommand{\SO}{\mathrm{SO}(3)}
\newcommand{\Inat}{I^{\natural}}
\newcommand{\td}{\widetilde d}
\newcommand{\tC}{\widetilde C}
\newcommand{\pill}{\mathcal{P}}

\begin{document}

\title{The S-complex differential of torus knots}

\author{Bernd Johannes Wuebben}
\subjclass[2020]{57K18 (primary); 57K10, 57R58, 11P21 (secondary)}
\keywords{Singular instanton homology, S-complexes, torus knots, traceless character varieties,
signature, Alexander polynomial, numerical semigroups, pillowcase}
\date{First version: October 2, 2026. This version: October 3, 2026}

\begin{abstract}
For a torus knot $K=T(p,q)$ the traceless $\SU$ character variety is finite and nondegenerate, so
Daemi and Scaduto's S-complex for singular instanton homology can be built on it. The resulting complex
has rank $1+|\sigma(K)|$, where $\sigma$ is the signature, and by theorems of Li and Ye and of Kronheimer
and Mrowka its homology $\Inat(K)$ has rank $\|\Delta_K\|_1$, the sum of the absolute values of the
coefficients of the Alexander polynomial; so the differential has rank $R(p,q)=\frac12(1+|\sigma|-\|\Delta\|_1)$.
We express $R(p,q)$ through floor sums and show that, for coprime $2\le p<q$, $R(p,q)\ge\lfloor(p-2b)^2/4\rfloor$,
where $b\in\{1,\dots,p-1\}$ is the inverse of $q$ modulo $p$, and that the differential vanishes if and
only if $p=2$, or $p$ is odd and $q\equiv2$ or $q\equiv p+2(-1)^{(p-1)/2}\pmod{2p}$; the `if' direction is
due to Daemi and Scaduto. By a result of Daemi and Scaduto, the component of the differential to the
reducible generator is nonzero over $\Q$ exactly when their invariant $h$ is positive. We show that the
components from degree $3$ to degree $2$ and from degree $1$ to degree $0$ have ranks $\lfloor R/2\rfloor$
and $\lceil R/2\rceil$, independently of the choices in the construction. In particular the
$\Z/4$-graded group $\Inat(T(p,q);\Q)$ depends only on $\sigma$ and $\|\Delta\|_1$, and we give it explicitly.
In each small-perturbation pillowcase computation of Hedden, Herald and Kirk, a bigon ends at the
reducible generator exactly when $h>0$. We formulate a chain-level conjecture identifying the two
complexes over $\mathbb F_2$; it predicts this pattern whenever the map to the reducible is nonzero
modulo $2$.
\end{abstract}

\maketitle

\section{Introduction}

\subsection{Background}
Kronheimer and Mrowka's reduced singular instanton homology $\Inat(K)$ of a knot $K\subset S^3$ is
built from traceless $\SU$ representations of $\pi_1(S^3\setminus K)$, those sending every meridian to
a conjugate of $\operatorname{diag}(i,-i)$~\cite{KM1,KM2}. Daemi and Scaduto refined the underlying
chain complex to an \emph{S-complex}~\cite[Def.~3.21]{DSEquiv}
\[
\tC_*(K)=C_*(K)\oplus C_{*-1}(K)\oplus\Z,\qquad
\td=\begin{pmatrix} d&0&0\\ v&-d&\delta_2\\ \delta_1&0&0\end{pmatrix}.
\]
Here $C_*(K)$ is generated by the irreducible critical points of the singular Chern--Simons functional,
and the summand $\Z$, in degree $0$, by the reducible (abelian) one. The maps $\delta_1$ and $\delta_2$
count instantons to and from the reducible, and the coefficients of $v\colon C_*\to C_{*-2}$ are the
degrees of holonomy maps to the circle on one-dimensional moduli spaces of instantons between
irreducibles~\cite[\S3.3, (3.15)]{DSEquiv}. Gradings are homological, differentials have degree $-1$, and degrees
are read modulo $4$. The S-complex structure is the map $\chi\colon\tC_*\to\tC_{*+1}$,
$\chi(\alpha,\beta,n)=(0,\alpha,0)$, which satisfies $\chi^2=0$ and
$\chi\td+\td\chi=0$~\cite[\S4.1]{DSEquiv}. A \emph{morphism} of S-complexes is a chain map of degree $0$
that commutes with $\chi$ and induces the identity on the reducible summand~\cite[Def.~3.28,
\S4.1]{DSEquiv}; an \emph{S-isomorphism} is an invertible morphism; and an \emph{S-chain homotopy
equivalence} is a pair of morphisms whose composites are chain homotopic to the identity through
homotopies that anticommute with $\chi$~\cite[Def.~3.30, \S4.1]{DSEquiv}. The S-chain homotopy type of
$\tC(K)$ depends only on the oriented knot $K$~\cite[Thm.~3.34, Rem.~3.35]{DSEquiv}, its homology is
$\Inat(K)$ up to a shift of grading~\cite[Thm.~8.9]{DSEquiv}, and its local equivalence class carries
Frøyshov-type concordance invariants, among them the integer $h(K)$~\cite[\S5]{DSEquiv}. Daemi and
Scaduto also define an equivariant homology $\widehat I(K)$, a module over $\Z[x]$~\cite[\S4.3,
Thm.~1.4]{DSEquiv}.

For a torus knot $K=T(p,q)$ the traceless character variety is finite and nondegenerate, with
$|\sigma(K)|/2$ irreducible points and one abelian point~\cite[\S5.1]{HHK1},~\cite[\S9.5]{DSEquiv}. The
S-complex can therefore be formed with a holonomy perturbation that is trivial near this critical set,
with one generator for each irreducible character and one for the abelian character~\cite[\S4.5, proof
of Thm.~4]{DScaduto}; we call such an S-complex \emph{built on the character variety}. The irreducible
generators lie in degrees $1$ and $3$~\cite[\S1, Cor.~3 and its proof]{DScaduto}, so $d=0$, $C_*(K)$ has
rank $|\sigma|/2$, and $\tC(K)$ has rank $1+|\sigma|$. Li and Ye showed that Kronheimer and Mrowka's
instanton knot homology $KHI(K)$ of a torus knot has dimension $\|\Delta_K\|_1$, the sum of the absolute
values of the coefficients of the Alexander polynomial~\cite[Cor.~1.8]{LiYe}; by Kronheimer and Mrowka's
isomorphism $\Inat(K;\Q)\cong KHI(K;\Q)$~\cite[Prop.~1.4]{KM2}, the homology of $\tC(K;\Q)$ has the same
rank. Hence the
rank of the S-complex differential of $T(p,q)$ is
\[
R(p,q):=\operatorname{rank}_\Q\td=\tfrac12\bigl(1+|\sigma(T(p,q))|-\|\Delta_{T(p,q)}\|_1\bigr)
\]
(Proposition~\ref{prop:rankformula}); this consequence is also recorded, with its values for $T(3,n)$,
in~\cite[Cor.~1.4]{WuebbenI}. This raises two questions. For which torus knots does the
differential vanish, so that the unperturbed singular Chern--Simons functional is perfect? And which
components of $\td$, the map $\delta_1$ to the reducible or the map $v$ between irreducibles, carry the
differential when it does not vanish?

Daemi and Scaduto answered the first question in one direction: they showed that $1+|\sigma|=\|\Delta\|_1$,
and hence $\td=0$ and $h=0$, for two explicit families of torus knots~\cite[(9.36), Prop.~9.37]{DSEquiv}. They also computed $T(3,4)$ and $T(3,5)$, where
$\delta_1\neq0$~\cite[\S\S9.3--9.4]{DSEquiv}. Hedden, Herald and Kirk computed the perturbed traceless
character varieties of several torus knots. Comparing Khovanov homology with the Alexander polynomial,
they determined the $\Z/4$-graded group $\Inat(T(3,n))$ for $n\le38$, whose rank is $1+|\sigma|$ or
$|\sigma|-1$, so that the differential has rank at most one~\cite[\S12.4]{HHK1}. They also noted that
Kronheimer and Mrowka's conjectured isomorphism between $KHI$ and knot Floer homology~\cite[Conj.~7.25]{KMsut}
would force the differential to have large rank in general~\cite[\S12.6]{HHK1}. With Li and Ye's theorem the rank of the differential is
an explicit function of $\sigma$ and $\Delta$, and Theorem~\ref{thm:main} below confirms this expectation.

\subsection{Results}
Throughout, $p,q\ge2$ are coprime, $T(p,q)=T(q,p)$ is the positive torus knot, and signatures are
normalized so that positive torus knots have $\sigma<0$; where needed we also allow $q=1$, for which
$T(p,1)$ is the unknot and $R(p,1)=0$. For a statement $\mathcal S$ we write $[\mathcal S]=1$ if $\mathcal S$ holds and $[\mathcal S]=0$
otherwise. For $p\ge2$ let $b=b_p(q)\in\{1,\dots,p-1\}$ be the inverse of $q$ modulo $p$, and put
\[
u_p(q)=p-2b_p(q),\qquad A_j=\Bigl\lfloor\frac{jq}{p}\Bigr\rfloor,\qquad F(n)=\sum_{j=0}^{n-1}A_j,
\qquad m(p,q)=\sum_{\substack{1\le k\le p-2\\ k\equiv p\ (\mathrm{mod}\ 2)}}\Bigl\lfloor\frac{kq}{2p}\Bigr\rfloor ;
\]
when the pair matters we write $F_{p,q}$. Since $T(p,q)=T(q,p)$, we have $R(p,q)=R(q,p)$, although the formula below is not symmetric in $p$
and $q$.

\begin{theorem}[Floor-sum formula and lower bound]\label{thm:main}
For coprime $p\ge2$ and $q\ge1$,
\[
R(p,q)=F(b)+F(p-b)-2\,m(p,q)\qquad\text{and}\qquad
R(p,q+2p)=R(p,q)+\Bigl\lfloor\frac{u_p(q)^2}{2}\Bigr\rfloor.
\]
If moreover $p<q$, then
\[
R(p,q)\ \ge\ \Bigl\lfloor\frac{u_p(q)^2}{4}\Bigr\rfloor.
\]
\end{theorem}

We call the second identity the \emph{step law}. The quantity $|u_p(q)|$ measures how far $b$ is from
$p/2$. In terms of the semigroup $\langle p,q\rangle$, $R$ compares the counting function of the
semigroup at two points with its value at their midpoint, and the bound says that $R$ is small only
when the two points are nearly symmetric (Remark~\ref{rem:semigroup}). The bound is sharp: for even $p\ge4$, $b_p(p+1)=1$
and $R(p,p+1)=((p-2)/2)^2=u_p(p+1)^2/4$. In particular $R$ is unbounded, as Hedden, Herald and Kirk
anticipated~\cite[\S12.6]{HHK1}. The smallest torus knot with $R\ge2$, by genus and by crossing number,
is $T(4,7)$, where $R=2$; for $T(5,11)$, $R=4$.

\begin{theorem}[Vanishing of the differential]\label{thm:perfect}
Let $2\le p<q$ be coprime. The S-complex differential of $T(p,q)$ vanishes, equivalently
$\operatorname{rank}_\Q\Inat(T(p,q))=1+|\sigma(T(p,q))|$, equivalently $\Delta_{T(p,q)}$ has exactly
$1+|\sigma(T(p,q))|$ nonzero coefficients, if and only if one of the following holds:
\begin{enumerate}
\item $p=2$;
\item $p$ is odd and $q\equiv2\pmod{2p}$;
\item $p\equiv1\pmod4$ and $q\equiv p+2\pmod{2p}$, or $p\equiv3\pmod4$ and $q\equiv p-2\pmod{2p}$.
\end{enumerate}
\end{theorem}

The torus knots in (2) and (3) are exactly the two families of~\cite[Prop.~9.37]{DSEquiv}, so the new
content of Theorem~\ref{thm:perfect} is the converse: no other torus knot has a perfect unperturbed
Chern--Simons functional. For $p=3$ it says that the differential of $T(3,n)$ vanishes exactly when
$n\equiv1,2\pmod6$, in agreement with the computations of Hedden, Herald and Kirk for
$n\le38$~\cite[\S12.4]{HHK1}. For torus knots $h(K)=\Upsilon_K(1)-\sigma(K)/2$~\cite[\S1, Prop.~1 and the
paragraph following it]{DScaduto}, where $\Upsilon_K$ is the invariant of Ozsv\'ath, Stipsicz and
Szab\'o~\cite{OSS}, normalized so that $\Upsilon_{T(2,3)}(1)=-1$. Jabuka and Van Cott showed that the knots
in (1)--(3) are exactly the torus knots with $\Upsilon_K(1)=\sigma(K)/2$ that become the unknot after a
single pinch move of Batson~\cite{Batson} (a nonorientable band move)~\cite[Prop.~1.6, Thms.~1.7--1.8]{JVC}.
Theorem~\ref{thm:perfect} can therefore be read as: the differential of a torus knot vanishes if and only
if $h=0$ and one pinch move unknots it.

For the second question, let $N_1=\lceil|\sigma|/4\rceil$ and $N_3=\lfloor|\sigma|/4\rfloor$. The complex
$C_*(K)$ is $C_1\oplus C_3$, with $\operatorname{rank}C_1=N_1$ and $\operatorname{rank}C_3=N_3$~\cite[\S1,
Cor.~3]{DScaduto}. Since $v$ has degree $-2$, it exchanges $C_1$ and $C_3$; write $v=v_{13}+v_{31}$ with
$v_{13}\colon C_1\to C_3$ and $v_{31}\colon C_3\to C_1$. In $\tC$ the group $C_1$ appears in degrees $1$ and
$2$, and $C_3$ in degrees $3$ and $0$; thus $v_{31}$ is the component of $\td$ from degree $3$ to degree
$2$, while $\delta_1$ and $v_{13}$ together form the component from degree $1$ to degree $0$. The map $v$
itself is not an invariant of $K$: an S-isomorphism can add to it a map of rank one
(Remark~\ref{rem:vnotinv}).

\begin{theorem}[Components of the differential]\label{thm:split}
Let $K=T(p,q)$, and use rational coefficients. For $w\in C_3$ write $w\,\delta_1$ for the map
$c\mapsto\delta_1(c)\,w$.
\begin{enumerate}
\item Any two S-complexes of $K$ built on its character variety are S-isomorphic. Every S-isomorphism
from $\tC$ to $\tC'$ is given by a pair $(P,w)$, with $P\colon C\to C'$ an isomorphism preserving degrees and
$w\in C'_3$, and then $v'=PvP^{-1}+w\,\delta_1P^{-1}$ and $\delta_1'=\delta_1P^{-1}$.
\item \cite[Prop.~4.15]{DSEquiv} $\delta_1\neq0$ if and only if $h(K)>0$.
\item $\operatorname{rank}v_{31}$ and $\operatorname{rank}(v_{13}|_{\ker\delta_1})$ do not depend on the
choice of S-complex built on the character variety, and
\[
\operatorname{rank}v_{31}=\Bigl\lfloor\frac{R(p,q)}2\Bigr\rfloor,\qquad
\operatorname{rank}(v_{13}|_{\ker\delta_1})=\Bigl\lceil\frac{R(p,q)}2\Bigr\rceil-[h(K)>0];
\]
equivalently, the components of $\td$ from degree $3$ to degree $2$ and from degree $1$ to degree $0$
have ranks $\lfloor R(p,q)/2\rfloor$ and $\lceil R(p,q)/2\rceil$.
\item If $h(K)=0$, then $\delta_1=0$, and $v$ is an endomorphism of $C_*(K;\Q)$ of rank $R(p,q)$, well
defined up to conjugation: it is the action of $x$ on the $\Q[x]$-torsion of Daemi and Scaduto's
equivariant homology $\widehat I(K;\Q)$.
\end{enumerate}
\end{theorem}

\begin{corollary}\label{cor:graded}
Let $K=T(p,q)$ and $R=R(p,q)$. In the grading of the homology of $\tC(K;\Q)$, in which the reducible
generator has degree $0$, $\Inat(K;\Q)$ has rank vector
\[
\bigl(N_3+1-\lceil R/2\rceil,\ N_1-\lceil R/2\rceil,\ N_1-\lfloor R/2\rfloor,\ N_3-\lfloor R/2\rfloor\bigr)
\]
in degrees $0,1,2,3$. The absolute grading of $\Inat$ used in~\cite[Thm.~8.9]{DSEquiv}, that
of~\cite[\S4.5]{KM2}, is obtained from this one by adding $\sigma(K)$ modulo $4$.
\end{corollary}

Earlier determinations of the $\Z/4$-graded group include $T(4,5)$~\cite[\S11]{KMfilt}, $T(3,n)$ with
$n\le38$~\cite[\S12.4]{HHK1}, and, since $\td=0$ there, the two families of~\cite[Prop.~9.37]{DSEquiv}.

Item (1) of Theorem~\ref{thm:split} follows from Daemi and Scaduto's definitions, item (2) is their
result, and the identification in (4) is formal. The new content of Theorem~\ref{thm:split} is (3),
together with its consequences: the rank of $v$ in (4) and Corollary~\ref{cor:graded}. It rests on
Lemma~\ref{lem:KM}. Kronheimer and Mrowka's operator $\mu(\bar\Sigma)$ on $KHI(K)$, whose generalized
eigenspaces are the summands in Alexander grading $j$, has degree $2$ modulo $4$, and these summands have
Euler characteristics equal to the coefficients of $\Delta_K$ up to sign~\cite{KMAlex,Lim}. With Li and Ye's
computation of their dimensions this gives $\{\beta_0-\beta_2,\beta_1-\beta_3\}=\{0,\pm1\}$, where $\beta_l$ is
the rank of $\Inat(K;\Q)$ in degree $l$. Since $\tC$ has ranks $(N_3+1,N_1,N_1,N_3)$, the ranks of the
components of degrees $3\to2$ and $1\to0$ then differ by at most one, and they sum to $R$. Kronheimer and
Mrowka gave this argument for $T(4,5)$~\cite[\S11]{KMfilt}.

Item (4) produces torus knots with $\delta_1=0$ and $v\neq0$ over $\Q$. Nonzero maps $v$ for torus knots
were known before in two settings. (i) Daemi and Scaduto also use local coefficients, over rings
containing a variable that they call $T$; when $T^4-1$ is invertible, the S-complex of a knot is S-chain
homotopy equivalent to that of the torus knot $T(2,|\sigma|+1)$ or of its mirror, according to the sign
of $\sigma$~\cite[\S1, Thm.~10]{DScaduto}, and over the field $\mathrm{Frac}((\Z/r)[T^{\pm1}])$, $r$ prime,
the S-complex of $T(p,q)$ has $\delta_1v^{|\sigma|/2-1}\neq0$~\cite[\S4.5, proof of Thm.~4]{DScaduto}, so
$v\neq0$ once $|\sigma|\ge4$. (ii) Over $\Z$, $\delta_1v\neq0$ whenever $h\ge2$~\cite[Prop.~4.15]{DSEquiv},
for instance for $T(5,6)$~\cite[\S4.1, Cor.~4]{DScaduto}. The smallest torus knot with $\delta_1=0$ and
$v\neq0$ over $\Q$, by genus and by crossing number, is $T(4,5)$, with $\sigma=-8$, $\|\Delta\|_1=7$, $R=1$
and $h=0$: its S-complex has $\delta_1=0$, $v_{31}=0$ and $v_{13}$ of rank one (Example~\ref{ex:45}). Other
examples are $T(4,9)$ ($R=2$), $T(6,7)$ ($R=4$) and $T(5,11)$ ($R=4$). For $T(3,n)$ the differential is
nonzero exactly when $n\equiv4,5\pmod6$ (Corollary~\ref{cor:3n}); then $h=1$, and up to S-isomorphism over $\Q$ the
differential is $\delta_1$ alone (Remark~\ref{rem:3n}).

\subsection{The pillowcase}
Hedden, Herald and Kirk computed pillowcase Floer homology for torus knots from a tangle decomposition
in which one tangle is trivial and carries the earring~\cite[\S\S10--11]{HHK2}. For small holonomy
perturbations the pillowcase complex, over $\mathbb F_2$, has $1+|\sigma|$ generators: a generator $r_+$,
the perturbation of the abelian character, in grading $\sigma\bmod4$, and two generators $x_i^-$ and
$x_i^+$ for each irreducible traceless character, with $\operatorname{gr}x_i^+=\operatorname{gr}x_i^-+1$
(labels as in~\cite[\S11]{HHK2}). Thus both complexes have one generator for the abelian character and
two, in adjacent gradings, for each irreducible character. If every bigon runs from some $x_i^-$ either
to $r_+$ or to some $x_j^+$ with $j\neq i$, then $\chi(x_i^-)=x_i^+$ and $\chi(x_i^+)=\chi(r_+)=0$ make the
pillowcase complex an S-complex $\pill(K)$ with reducible generator $r_+$, in which $x_i^-$ and $x_i^+$ play
the roles of the two copies of an irreducible generator. Comparing the computations of Hedden, Herald and
Kirk with Theorem~\ref{thm:split} suggests the following.

\begin{conjecture}\label{conj:pillow}
Let $K=T(p,q)$, and consider the pillowcase complex of~\cite[\S\S10--11]{HHK2} over $\mathbb F_2$, for the
decomposition in which the earring lies on the trivial tangle and the integers $r,s$ with $pr+qs=1$ satisfy
$0<r<q$, with zero bounding cochains and a holonomy perturbation small enough that the complex has exactly
the $1+|\sigma|$ generators $r_+$, $x_i^\pm$ described above.
\begin{enumerate}
\item[(a)] Every bigon runs from some $x_i^-$ either to $r_+$ or to some $x_j^+$ with $j\neq i$.
\item[(b)] The S-complex $\pill(K)$, with $\sigma$ subtracted from its grading, is S-chain homotopy
equivalent to $\tC(K;\mathbb F_2)$.
\end{enumerate}
\end{conjecture}

The decomposition of~\cite[\S10]{HHK2} depends on the choice of $(r,s)$. The condition $0<r<q$ determines it,
since $pr\equiv1\pmod q$, and all the choices printed in~\cite[\S11]{HHK2} satisfy it. The restrictions are needed.
For larger perturbations the complex of $T(3,4)$ becomes perfect, and with the earring on the other tangle
the homology of $T(4,5)$ has the wrong rank (Section~\ref{sec:pillow}). For $T(3,4)$ with $(r,s)=(7,-5)$, a
numerical computation gives a complex with zero bounding cochains, no bigon, and homology of rank $7$ rather
than $5$~\cite{Wpill}.

The S-complex $\tC(K;\mathbb F_2)$ has $d=0$, and under (a) so does $\pill(K)$. For an S-complex with
$d=0$, the map to the reducible is nonzero if and only if Daemi and Scaduto's invariant $h$ is
positive~\cite[Prop.~4.15, Rem.~4.2]{DSEquiv}, and $h$ depends only on the S-chain homotopy type. Hence the
conjecture implies that some $x_i^-$ is joined to $r_+$ by an odd number of bigons if and only if
$\delta_1\otimes\mathbb F_2\neq0$, that is, if and only if the homomorphism $\delta_1\colon C_1(K;\Z)\to\Z$
takes an odd value. This condition implies $h(K)>0$. For $T(3,4)$ and $T(3,5)$ it is equivalent to $h(K)>0$. There, for each
irreducible generator $\alpha$ of degree $1$, the moduli space defining $\delta_1(\alpha)$ consists of a single
instanton~\cite[\S9.3, and the same argument in \S9.4]{DSEquiv}, so $\delta_1(\alpha)=\pm1$.
Section~\ref{sec:pillow} lists the evidence: every torus knot with nonzero pillowcase differential
in~\cite{HHK2} is consistent with (a) and with the components predicted by Theorem~\ref{thm:split}, including $T(4,5)$, whose bigon avoids $r_+$ as Hedden, Herald and
Kirk point out~\cite[\S11.8]{HHK2}, and $T(4,7)$, with one bigon into $r_+$ and one between irreducible
generators, exactly as Theorem~\ref{thm:split} requires. For the torus knots $T(3,n)$, part~(a) holds, and the
pillowcase complex has the graded generators of $\tC(K)$ and a differential of rank $R$ consisting of a single
component of type $\delta_1$, and combined with a cobordism inequality of Daemi and Scaduto this gives part~(b) for every
$T(3,n)$~\cite[Thm.~1.3]{Wpill}. Beyond
torus knots, for the pretzel knots $P(-2,3,q)$ the pillowcase complex of a decomposition obtained from that
of $T(3,5)$ by Dehn twists along the Conway sphere is again an S-complex over $\mathbb F_2$, but its only nonzero
component is the differential $d$ of the irreducible part, which is nonzero for $q\ge9$~\cite[Rem.~4.20]{WuebbenII}.

\subsection{Outline}
The formula in Theorem~\ref{thm:main} combines Moree's count $\|\Delta_{T(p,q)}\|_1=2ab-1$, where
$ap+bq=pq+1$, with a reorganization of Litherland's lattice-point formula for the signature
(Sections~\ref{sec:prelim} and~\ref{sec:formula}). From the formula we derive the step law and a
reflection law relating $R(p,q)$ and $R(p,2p-q)$; via Moree's count these are the recursions of Gordon,
Litherland and Murasugi for the signature (Section~\ref{sec:laws}). The lower bound is proved in
Section~\ref{sec:bound}. If $p$ or $q$ is even, it follows from a telescoping expression for $R$ and a
reciprocity law relating $u_p(q)$ and $u_q(p)$. If both are odd, the step and reflection laws express
$R(p,q)$ through $R(p_1,p)$ with $p_1<p$, one step of a Euclidean algorithm with even quotients, and
reciprocity controls how $|u|$ changes under this step; because the reflection law enters with a minus
sign, the induction carries an upper bound along with the lower bound. Theorem~\ref{thm:perfect} follows
in Section~\ref{sec:perfect}: the bound gives $R\ge1$ unless $|u_p(q)|\le1$, and then the step law reduces
$R(p,q)$ to four residues of $q$ modulo $2p$. Section~\ref{sec:split} proves Theorem~\ref{thm:split}, and
Section~\ref{sec:pillow} compares it with the pillowcase computations.

\section{Preliminaries}\label{sec:prelim}

\subsection{The S-complex of a torus knot}

\begin{proposition}\label{prop:rankformula}
Let $K$ be a torus knot. The S-complex of $K$ can be built on its traceless character variety, with
$d=0$, $\delta_2=0$, and $C_*(K)=C_1\oplus C_3$, $\operatorname{rank}C_1=N_1$, $\operatorname{rank}C_3=N_3$.
Over $\Q$ its differential has rank $R=\frac12(1+|\sigma(K)|-\|\Delta_K\|_1)$, independently of the choices
made in the construction.
\end{proposition}

\begin{proof}
The traceless character variety of $K$ consists of $|\sigma|/2$ irreducible points and one abelian
point, all nondegenerate~\cite[\S9.5]{DSEquiv}, and the complex $C_*(K)$ built on it has zero
differential and the stated gradings~\cite[\S1, Cor.~3 and its proof]{DScaduto}. Since $C_2=0$, the map
$\delta_2\colon\Z\to C_{-2}=C_2$ vanishes. The homology of $\tC$ is isomorphic to $\Inat(K)$ up to a
grading shift~\cite[Thm.~8.9]{DSEquiv}, and $\operatorname{rank}\Inat(K)=\dim KHI(K)=\|\Delta_K\|_1$
by~\cite[Prop.~1.4]{KM2} and~\cite[Cor.~1.8]{LiYe}. Since $\dim H=\dim\tC-2\operatorname{rank}\td$ and
$\dim\tC=1+|\sigma|$, the formula follows; both sides depend only on $K$.
\end{proof}

\subsection{Signature and Alexander polynomial}
Let $\Gamma=\langle p,q\rangle=\{ip+jq:i,j\in\Z_{\ge0}\}$. Up to a unit,
$\Delta_{T(p,q)}(t)=(1-t)(1-t^{pq})/((1-t^p)(1-t^q))$, a classical formula (see
e.g.~\cite[\S9.6]{DSEquiv}), and this equals $(1-t)\sum_{\gamma\in\Gamma}t^\gamma$~\cite[Thm.~1]{Moree}, so
every coefficient lies in $\{-1,0,1\}$. Let $a\ge1$ be defined by $ap+bq=pq+1$; equivalently $a$ is the
inverse of $p$ modulo $q$.

\begin{lemma}[Moree~{\cite[Lemma~1, Cor.~1]{Moree}}]\label{lem:moree}
$\|\Delta_{T(p,q)}\|_1=2ab-1$.
\end{lemma}

For primes $p$ and $q$ this count is due to Carlitz; see~\cite{Moree}. The signature is given by
Litherland's lattice-point count~\cite{Litherland},~\cite[Prop.~1.1]{BO}, where $(i,j)$ ranges over
$1\le i\le p-1$ and $1\le j\le q-1$:
\[
\sigma(T(p,q))=\#\bigl\{(i,j)\colon \tfrac ip+\tfrac jq\notin(\tfrac12,\tfrac32)\bigr\}
-\#\bigl\{(i,j)\colon \tfrac ip+\tfrac jq\in(\tfrac12,\tfrac32)\bigr\},
\]

\begin{lemma}\label{lem:sigma}
$|\sigma(T(p,q))|=(p-1)(q-1)-4m(p,q)$.
\end{lemma}

\begin{proof}
No sum $i/p+j/q$ equals $\frac12$ or $\frac32$: from $2iq+2jp\in\{pq,3pq\}$ we get $p\mid 2i$, hence
$2i=p$ and $j\in\{0,q\}$, which is excluded. The involution $(i,j)\mapsto(p-i,q-j)$ exchanges the sums
below $\frac12$ with those above $\frac32$, so the number of sums outside $(\frac12,\frac32)$ is twice the
number $m$ of sums below $\frac12$. For $(i,j)$ with $i/p+j/q<\frac12$, both $(p-i,j)$ and $(i,q-j)$ give
sums in $(\frac12,\frac32)$, and these two injections have disjoint images; hence $\sigma\le0$ and
$|\sigma|=(p-1)(q-1)-4m$. Finally, the condition $2(iq+jp)<pq$ forces $2i<p$, and writing $k=p-2i$ it
reads $j<kq/(2p)$; since $p\nmid kq$, the number of such $j\ge1$ is $\lfloor kq/2p\rfloor$, and $k$ runs
over $1\le k\le p-2$ with $k\equiv p\pmod2$. So $m=m(p,q)$.
\end{proof}

\section{The floor-sum formula}\label{sec:formula}

We combine Lemmas~\ref{lem:moree} and~\ref{lem:sigma} with a count of the floor sums $F$.

\begin{proof}[Proof of the formula in Theorem~\ref{thm:main}]
Let $g=(p-1)(q-1)/2$, the genus of $T(p,q)$. The residues of $jq$ modulo $p$ for $0\le j\le p-1$ are
$0,\dots,p-1$, so $F(p)=\frac1p\bigl(q\binom p2-\binom p2\bigr)=g$. For $0\le i\le b-1$ we have
$(p-b+i)q=ap+iq-1$, so $A_{p-b+i}=a-1$ for $i=0$ and $A_{p-b+i}=a+A_i$ for $1\le i\le b-1$ (as
$p\nmid iq$). Summing, $F(p)-F(p-b)=F(b)+ab-1$, that is, $F(b)+F(p-b)=g+1-ab$. By
Lemmas~\ref{lem:moree} and~\ref{lem:sigma},
\[
R=\tfrac12\bigl(1+2g-4m-(2ab-1)\bigr)=1+g-ab-2m=F(b)+F(p-b)-2m .\qedhere
\]
\end{proof}

\begin{remark}\label{rem:semigroup}
In terms of the semigroup, $F(p-b)$, $F(b)$ and $m$ are the numbers of elements of $\Gamma$ in
$[0,X_1]$, $[0,X_2]$ and $[0,M)$, where $X_1=pq-p-q-bq$, $X_2=bq-p-q-1$ and $M=\frac12(pq-2p-2q)$. Thus
\[
R=\#(\Gamma\cap[0,X_1])+\#(\Gamma\cap[0,X_2])-2\,\#(\Gamma\cap[0,M)),\qquad X_1+X_2=2M-1,
\]
so $R$ compares the counting function of $\Gamma$ at two points with its value at their midpoint. For
$p<q$, Theorem~\ref{thm:main} shows that $R$ vanishes only when $b$ is close to $p/2$, that is, when
$X_1$ and $X_2$ are nearly symmetric about $M$.
\end{remark}

\section{The step and reflection laws}\label{sec:laws}

Write $u=u_p(q)$; it has the parity of $p$. The following two recursions are, via
Lemma~\ref{lem:moree}, the recursions of Gordon, Litherland and Murasugi for the
signature~\cite[Thm.~5.2]{GLM}; we derive them from the floor-sum formula.

\begin{lemma}[Step law]\label{lem:step}
For $p\ge2$ and $q\ge1$ coprime, $R(p,q+2p)=R(p,q)+\lfloor u^2/2\rfloor$.
\end{lemma}

\begin{proof}
The inverse $b$ is unchanged. Each $A_j$ increases by $2j$, so $F(n)$ increases by $n(n-1)$, and each
term $\lfloor kq/2p\rfloor$ increases by $k$, so $m$ increases by $\Sigma_1=\sum k$ over
$1\le k\le p-2$, $k\equiv p\pmod2$. Using $b^2+(p-b)^2=\frac12(p^2+u^2)$, the increase of $R$ is
$\frac12(p^2+u^2)-p-2\Sigma_1$. For $p$ odd, $\Sigma_1=\frac14(p-1)^2$ and the increase is
$\frac12(u^2-1)$; for $p$ even, $\Sigma_1=\frac p2(\frac p2-1)$ and the increase is $\frac12u^2$.
\end{proof}

\begin{lemma}[Reflection law]\label{lem:refl}
For coprime $p<q<2p$, $R(p,q)+R(p,2p-q)=\lfloor u^2/2\rfloor+[p\text{ odd}]$.
\end{lemma}

\begin{proof}
Put $q'=2p-q$; then $b_p(q')=p-b$. For $1\le j\le p-1$, $jq/p\notin\Z$, so
$\lfloor jq/p\rfloor+\lfloor jq'/p\rfloor=2j-1$, and the floor sums satisfy
$F_{p,q}(n)+F_{p,q'}(n)=(n-1)^2$ for $1\le n\le p$. Likewise
$\lfloor kq/2p\rfloor+\lfloor kq'/2p\rfloor=k-1$, so $m(p,q)+m(p,q')=\Sigma_2=\sum(k-1)$ over the same
$k$. Hence the sum is $(b-1)^2+(p-b-1)^2-2\Sigma_2=\frac12\bigl((p-2)^2+u^2\bigr)-2\Sigma_2$, which
equals $\frac12(u^2+1)$ for $p$ odd ($\Sigma_2=\frac14(p-1)(p-3)$) and $\frac12u^2$ for $p$ even
($\Sigma_2=(\frac p2-1)^2$).
\end{proof}

\section{The lower bound}\label{sec:bound}

\subsection{\texorpdfstring{$p$ even}{p even}}

\begin{proposition}\label{prop:even}
Let $p\ge2$ be even and $q\ge1$ odd, and let $s=|u_p(q)|/2$. Then
$R(p,q)=\sum_{i=0}^{s-1}\bigl(A_{p/2+i}-A_{p/2-s+i}\bigr)\ge s\lfloor sq/p\rfloor$. In particular
$R(p,q)\ge u_p(q)^2/4$ when $q>p$.
\end{proposition}

\begin{proof}
For $p$ even the indices $k=2l$ give $m(p,q)=\sum_{l=1}^{p/2-1}\lfloor lq/p\rfloor=F(p/2)$. The
expression $F(b)+F(p-b)$ is symmetric under $b\leftrightarrow p-b$, and $\{b,p-b\}=\{p/2-s,p/2+s\}$.
Hence $R=[F(p/2+s)-F(p/2)]-[F(p/2)-F(p/2-s)]$, which telescopes to the stated sum. Each term is at
least $\lfloor sq/p\rfloor$, because $\lfloor x+y\rfloor-\lfloor x\rfloor\ge\lfloor y\rfloor$. If
$q>p$ then $\lfloor sq/p\rfloor\ge s$.
\end{proof}

\subsection{Reciprocity}
We extend $u$ by $u_1(q)=-1$, which is $p-2b$ with $p=b=1$.

\begin{lemma}\label{lem:recip}
For coprime $n_1,n_2\ge1$, $n_1u_{n_2}(n_1)+n_2u_{n_1}(n_2)=-2$. Consequently, if $n_2\ge3$ and $n_1\neq2$,
then $u_{n_2}(n_1)$ and $u_{n_1}(n_2)$ have opposite signs.
\end{lemma}

\begin{proof}
For $n_1,n_2\ge2$ let $c=b_{n_2}(n_1)$ and $e=b_{n_1}(n_2)$. Then $cn_1+en_2\equiv1$ modulo $n_1$ and modulo
$n_2$, and $n_1+n_2\le cn_1+en_2\le2n_1n_2-n_1-n_2$, so $cn_1+en_2=n_1n_2+1$ (the identity $ap+bq=pq+1$ of Section~\ref{sec:prelim} for the pair
$(n_2,n_1)$) and
$n_1(n_2-2c)+n_2(n_1-2e)=-2$. If $n_1=1$, then $u_{n_2}(1)=n_2-2$ and $1\cdot(n_2-2)+n_2\cdot(-1)=-2$. For
the sign: if $n_1\neq2$ then $u_{n_1}(n_2)\neq0$, since $b_{n_1}(n_2)=n_1/2$ would give
$n_1n_2/2\equiv1\pmod{n_1}$, hence $n_1/2\mid1$. If $u_{n_1}(n_2)\ge1$ then $n_1u_{n_2}(n_1)\le-n_2-2<0$, and if
$u_{n_1}(n_2)\le-1$ then $n_1u_{n_2}(n_1)\ge n_2-2>0$.
\end{proof}

\subsection{\texorpdfstring{$p$ odd, $q$ even}{p odd, q even}}

\begin{proposition}\label{prop:oddeven}
Let $p\ge3$ be odd and $q>p$ even. Then $R(p,q)\ge(u_p(q)^2-1)/4$.
\end{proposition}

\begin{proof}
Put $\ell=|u_p(q)|$ and $s'=|u_q(p)|/2$. Proposition~\ref{prop:even} applied to $(q,p)$ gives
$R=R(q,p)\ge s'\lfloor s'p/q\rfloor\ge0$, which settles the case $\ell=1$. Let $\ell\ge3$. By
Lemma~\ref{lem:recip}, $2ps'=|qu_p(q)+2|\ge q\ell-2$. Since $q\ge p+1$, this gives
$s'\ge\ell/2+(\ell-2)/(2p)>\ell/2$, so $s'\ge(\ell+1)/2$; and $s'p/q\ge\ell/2-1/q>(\ell-1)/2$, so
$\lfloor s'p/q\rfloor\ge(\ell-1)/2$. Hence $R\ge(\ell+1)(\ell-1)/4$.
\end{proof}

\subsection{\texorpdfstring{$p$ and $q$ odd}{p and q odd}}
This case is handled by a descent $(p,q)\mapsto(p_1,p)$, one step of a Euclidean algorithm with even
quotients.

\begin{lemma}\label{lem:euclid}
Let $3\le p<q$ be odd and coprime, and let $\rho\in(0,2p)$ be the residue of $q$ modulo $2p$; then
$\rho\neq p$ because $p\nmid q$. Put $(p_1,\varepsilon)=(\rho,+1)$ if $\rho<p$ and
$(p_1,\varepsilon)=(2p-\rho,-1)$ if $\rho>p$, and let $k\ge1$ be such that $q=2kp+\varepsilon p_1$. Write
$e_0=u_p(p_1)$, $e_1=u_{p_1}(p)$, $\theta_0=|e_0|$ and $\theta_1=|e_1|$. Then $p_1$ is odd and coprime to $p$,
$1\le p_1\le p-2$, and:
\begin{enumerate}
\item $|u_p(q)|=\theta_0$ and $|u_q(p)|=2k\theta_0+\varepsilon\theta_1$;
\item $\theta_1\le\theta_0$;
\item $R(p,q)=R(p_1,p)+k(\theta_0^2-1)/2$ if $\varepsilon=+1$, and
$R(p,q)=k(\theta_0^2-1)/2+1-R(p_1,p)$ if $\varepsilon=-1$.
\end{enumerate}
For example, $(p,q)=(5,13)$ gives $\rho=3$, $p_1=3$, $\varepsilon=+1$ and $k=1$.
\end{lemma}

\begin{proof}
Since $q\equiv\varepsilon p_1\pmod p$ and $u_p(-n)=-u_p(n)$, we have $u_p(q)=\varepsilon e_0$.
Lemma~\ref{lem:recip} for $(p_1,p)$ gives $p_1e_0+pe_1=-2$, so $e_0,e_1$ are odd of opposite signs, and
$p_1\theta_0\ge p\theta_1-2\ge(p_1+2)\theta_1-2\ge p_1\theta_1$, which is (2). Lemma~\ref{lem:recip} for
$(p,q)$ gives $pu_q(p)=-2-q\varepsilon e_0=-2k\varepsilon pe_0+pe_1$ (substitute $q=2kp+\varepsilon p_1$ and
$p_1e_0=-2-pe_1$), so $u_q(p)=-2k\varepsilon e_0+e_1$; the
two terms have the same sign if $\varepsilon=+1$ and opposite signs with $2k\theta_0>\theta_1$ if
$\varepsilon=-1$, which gives (1). For (3), $|u_p(n)|=\theta_0$ for $n\equiv\pm p_1\pmod p$. If $\varepsilon=+1$
apply Lemma~\ref{lem:step} $k$ times starting from $R(p,p_1)=R(p_1,p)$. If $\varepsilon=-1$ apply it $k-1$
times starting from $R(p,2p-p_1)$, which equals $(\theta_0^2+1)/2-R(p_1,p)$ by Lemma~\ref{lem:refl}.
\end{proof}

Because the reflection law enters with a minus sign, the induction must carry an upper bound alongside
the lower one.

\begin{proposition}\label{prop:oddodd}
For odd coprime $1\le p<q$,
\[
\frac{u_p(q)^2-1}{4}\ \le\ R(p,q)\ \le\ \frac{u_q(p)^2+3}{4}.
\]
\end{proposition}

\begin{proof}
Induction on $p$. For $p=1$, $R(1,q)=0$ and $u_1(q)^2=1$. For $p\ge3$, take the data of
Lemma~\ref{lem:euclid}; by induction $(\theta_1^2-1)/4\le R(p_1,p)\le(\theta_0^2+3)/4$.

Lower bound: if $\varepsilon=+1$, $R(p,q)\ge k(\theta_0^2-1)/2\ge(\theta_0^2-1)/4$. If $\varepsilon=-1$,
$R(p,q)\ge k(\theta_0^2-1)/2+1-(\theta_0^2+3)/4=(2k-1)(\theta_0^2-1)/4\ge(\theta_0^2-1)/4$.

Upper bound, with $\theta=|u_q(p)|$: if $\varepsilon=+1$, then $4R(p,q)\le\theta_0^2+3+2k(\theta_0^2-1)$, while
$\theta^2\ge4k^2\theta_0^2\ge(2k+1)\theta_0^2>\theta_0^2+2k(\theta_0^2-1)$. If $\varepsilon=-1$, then
$4R(p,q)\le2k(\theta_0^2-1)+4-(\theta_1^2-1)$ and $\theta=2k\theta_0-\theta_1$, and
expanding,
\[
\theta^2+3-\bigl[2k(\theta_0^2-1)+5-\theta_1^2\bigr]
=2(\theta_0-\theta_1)^2+2(2k-1)(k-1)\theta_0^2+4(k-1)\theta_0(\theta_0-\theta_1)+2(k-1)\ge0.\qedhere
\]
\end{proof}

\begin{proof}[Proof of the lower bound in Theorem~\ref{thm:main}]
If $p$ is even, this is Proposition~\ref{prop:even}. If $p$ is odd then $u_p(q)$ is odd and
$\lfloor u_p(q)^2/4\rfloor=(u_p(q)^2-1)/4$; use Proposition~\ref{prop:oddeven} if $q$ is even and
Proposition~\ref{prop:oddodd} if $q$ is odd. The step law is Lemma~\ref{lem:step}.
\end{proof}

\section{Perfectness}\label{sec:perfect}

\begin{lemma}\label{lem:tau01}
For coprime $p\ge2$ and $q$, $|u_p(q)|\le1$ if and only if $p=2$ or $q\equiv\pm2\pmod p$.
\end{lemma}

\begin{proof}
$u_p(q)=0$ means $p=2b$, which with $\gcd(b,p)=1$ forces $p=2$. $|u_p(q)|=1$ means $p$ odd and
$b=(p\pm1)/2\equiv\pm2^{-1}\pmod p$, that is, $q\equiv\pm2\pmod p$.
\end{proof}

\begin{lemma}\label{lem:special}
$R(2,n)=0$ for every odd $n\ge1$. For odd $p\ge3$, $R(p,2p-2)=1$ and
$R(p,p+2)=[p\equiv3\pmod4]$; moreover $R(p,p-2)=[p\equiv1\pmod4]$, including $R(3,1)=0$.
\end{lemma}

\begin{proof}
For $(2,n)$ the formula of Theorem~\ref{thm:main} gives $R(2,n)=0$, since $b=1$, $F(1)=0$ and
$m(2,n)=0$; in particular $R(p,2)=R(2,p)=0$. Since $b_p(2)=(p+1)/2$ and $b_p(-2)=(p-1)/2$,
$|u_p(n)|=1$ for $n\equiv\pm2\pmod p$, and Lemma~\ref{lem:refl} gives $R(p,2p-2)=1-R(p,2)=1$ and
$R(p,p+2)=1-R(p,p-2)=1-R(p-2,p)$. With $R(1,3)=0$ ($T(1,3)$ is the unknot), induction gives
$R(p,p+2)=0$ for $p\equiv1$ and $1$ for $p\equiv3\pmod4$, and then $R(p,p-2)=1-R(p,p+2)$.
\end{proof}

\begin{proof}[Proof of Theorem~\ref{thm:perfect}]
If $p=2$ then $R=0$ by Lemma~\ref{lem:special}. Let $p\ge3$. If $|u_p(q)|\ge2$, Theorem~\ref{thm:main}
gives $R\ge1$. So assume $|u_p(q)|=1$, that is, $p$ odd and $q\equiv\pm2\pmod p$
(Lemma~\ref{lem:tau01}). By the step law with $\lfloor u^2/2\rfloor=0$, $R(p,q)=R(p,q_0)$ for the residue
$q_0\in\{2,p-2,p+2,2p-2\}$ of $q$ modulo $2p$. Lemma~\ref{lem:special} gives:
$R=0$ for $q_0=2$; $R=1$ for $q_0=2p-2$; $R=0$ for $q_0=p+2$ iff $p\equiv1\pmod4$; and $R=0$ for $q_0=p-2$
iff $p\equiv3\pmod4$. These are the conditions stated.
\end{proof}

\begin{corollary}\label{cor:3n}
For $n\ge1$ coprime to $3$, $R(3,n)=0$ for $n\equiv1,2\pmod6$ and $R(3,n)=1$ for $n\equiv4,5\pmod6$.
\end{corollary}

\begin{proof}
For $p=3$, $|u_3(n)|=1$ for every $n$, so by the step law $R(3,n)$ depends only on $n$ modulo $6$, and
the four residues are covered by Lemma~\ref{lem:special}. For $n=4,5$ this
is~\cite[\S\S9.3--9.4]{DSEquiv}, and for $n\equiv1,2\pmod6$ it is~\cite[Prop.~9.37]{DSEquiv}.
\end{proof}

\section{Components of the differential}\label{sec:split}

\begin{lemma}[after Kronheimer and Mrowka~{\cite[\S11]{KMfilt}}]\label{lem:KM}
Let $K$ be a torus knot, and let $\beta_l$ be the rank of $\Inat(K;\Q)$ in degree $l\in\Z/4$. Then one of
$\beta_0-\beta_2$ and $\beta_1-\beta_3$ is $0$ and the other is $\pm1$. The same holds for the ranks of the
homology of $\tC(K;\Q)$, graded as in Section~\ref{sec:prelim}.
\end{lemma}

\begin{proof}
Let $\bar\Sigma$ be the closed surface obtained by capping off a Seifert surface of $K$, of genus $g$, in
Kronheimer and Mrowka's closure of the knot complement, and let $\mu=\mu(\bar\Sigma)$ be the corresponding
operator on $KHI(K;\Q)$. Its generalized eigenspaces $KHI(K,j)$ belong to the eigenvalues $2j$,
$|j|\le g$~\cite[Def.~2.5]{KMAlex}; $KHI(K,j)$ is the summand in Alexander grading $j$. With respect to
the canonical $\Z/2$ grading, the Euler characteristic of $KHI(K,j)$ is $-c_j$, where $c_j$ is the
coefficient of $t^j$ in the symmetrized Alexander polynomial~\cite[Thm.~1.1]{KMAlex},~\cite{Lim}. We give
$KHI(K;\Q)$ the relative $\Z/4$ grading carried over from $\Inat(K;\Q)$ along Kronheimer and Mrowka's
isomorphism; they state that this isomorphism respects the mod $4$ grading and that $\mu$ has degree $2$
for it~\cite[\S11]{KMfilt}. Only the relative grading is used, since the conclusion is unchanged by a
cyclic shift or a reversal of the grading.
\begin{enumerate}
\item[(i)] For torus knots $\dim KHI(K,j)=|c_j|$~\cite[Cor.~1.8]{LiYe}.
\item[(ii)] The central coefficient is $c_0=\pm1$. Indeed, up to a unit $\Delta_K=(1-t)\sum_{\gamma\in\Gamma}t^\gamma$
with $\Gamma=\langle p,q\rangle$ (Section~\ref{sec:prelim}), a polynomial of degree $2g$ whose coefficient of
$t^g$ is the central coefficient; so $c_0=\pm\bigl([g\in\Gamma]-[g-1\in\Gamma]\bigr)$, and exactly one of $g$
and $g-1$ lies in $\Gamma$, because $n\in\Gamma$ if and only if $2g-1-n\notin\Gamma$.
\item[(iii)] For $j\neq0$, $V_j=KHI(K,j)\oplus KHI(K,-j)$ is the generalized eigenspace of $\mu^2$ for the
eigenvalue $4j^2$. Since $\mu^2$ has degree $0$, $V_j$ is a graded subspace, and $\mu$ restricts to an
isomorphism of $V_j$ of degree $2$; so the ranks of $V_j$ in degrees $l$ and $l+2$ agree. Likewise
$KHI(K,0)$, the generalized kernel of $\mu^2$, is graded, and by (i) and (ii) it is one-dimensional.
\item[(iv)] Summing over $j$, the ranks of $KHI(K;\Q)$ in degrees $l$ and $l+2$ agree for the pair of
degrees not containing the degree of $KHI(K,0)$, and differ by one for the other pair.
\end{enumerate}
By the choice of grading the same holds for $\Inat(K;\Q)$, and by~\cite[Thm.~8.9]{DSEquiv}, which identifies
the homology of $\tC(K;\Q)$ with $\Inat(K;\Q)$ up to a shift of grading, it holds for the homology of
$\tC(K;\Q)$.
\end{proof}

The proof uses only that $\dim KHI(K;\Q)=\|\Delta_K\|_1$ and that the central coefficient of $\Delta_K$ is
$\pm1$.

\begin{proof}[Proof of Theorem~\ref{thm:split}]
(1) Let $\tC$ and $\tC'$ be S-complexes of $K$ built on its character variety. They are S-chain homotopy
equivalent~\cite[Thm.~3.34]{DSEquiv}. A morphism $\tC\to\tC'$ has the block form
\[
\begin{pmatrix}P&0&0\\ N&P&E_2\\ E_1&0&1\end{pmatrix},
\]
subject to the chain-map relations of~\cite[Def.~3.28]{DSEquiv}; the form is forced by commuting with
$\chi$~\cite[\S4.1]{DSEquiv}. The gradings force $N=0$ (it would map $C_1\to C'_0$ and
$C_3\to C'_2$, both zero) and $E_1=0$ (it would map $C_0=0$ to $\Q$), and the relations reduce to
$\delta_1=\delta_1'P$ and $v'P=Pv+E_2\delta_1$, where $E_2\in C'_3$ is the image of $1\in\Q$. An S-chain
homotopy has degree $1$, and $C$ is concentrated in degrees $1$ and $3$, so its diagonal blocks vanish.
Hence the diagonal blocks of a morphism and of its homotopy inverse are inverse to each other, and $P$
is invertible. With $w=E_2$ the
relations become $v'=PvP^{-1}+w\,\delta_1P^{-1}$ and $\delta_1'=\delta_1P^{-1}$. Conversely every such pair
$(P,w)$ defines an S-isomorphism.

(2) Since $d=0$, every element of $C_*$ is a cycle, and $h(K)>0$ if and only if $\delta_1\neq0$
by~\cite[Prop.~4.15]{DSEquiv}.

(3) The map $w\,\delta_1$ sends $C_1$ to $C_3$ and vanishes on $\ker\delta_1$, so it changes neither $v_{31}$
nor $v_{13}|_{\ker\delta_1}$, and $P$ conjugates both; hence their ranks are invariant. With $d=0$ and
$\delta_2=0$, $\td$ is determined by the map $(v,\delta_1)\colon C\to C\oplus\Q$, whose rank is
$[\delta_1\neq0]+\operatorname{rank}(v|_{\ker\delta_1})$, and
$\operatorname{rank}(v|_{\ker\delta_1})=\operatorname{rank}v_{31}+\operatorname{rank}(v_{13}|_{\ker\delta_1})$
because $C_3\subset\ker\delta_1$ and the two pieces land in the complementary summands $C_1$ and $C_3$.
Let $R_3=\operatorname{rank}v_{31}$, and let $R_1$ be the rank of
$(\delta_1,v_{13})\colon C_1\to\Q\oplus C_3$; these are the ranks of the components of $\td$ of degrees
$3\to2$ and $1\to0$. Then $R_1+R_3=R$ and
$R_1=[h>0]+\operatorname{rank}(v_{13}|_{\ker\delta_1})$ by (2). The complex $\tC$ has rank vector
$(N_3+1,N_1,N_1,N_3)$ in degrees $0,1,2,3$, so its homology has rank vector
\[
\beta=(N_3+1-R_1,\ N_1-R_1,\ N_1-R_3,\ N_3-R_3).
\]
Put $D=R_3-R_1$. Then $\beta_0-\beta_2=1-(N_1-N_3)+D$ and $\beta_1-\beta_3=(N_1-N_3)+D$. Since $|\sigma|$ is
even, $N_1-N_3\in\{0,1\}$, so these two differences are $D$ and $D+1$ in some order. By
Lemma~\ref{lem:KM} one of them is $0$ and the other is $\pm1$, so $D\in\{0,-1\}$. With $R_1+R_3=R$ this
gives $R_3=\lfloor R/2\rfloor$ and $R_1=\lceil R/2\rceil$.

(4) If $h=0$ then $\delta_1=0$ by (2), the transformation in (1) is conjugation by $P$, and $v$ has rank
$R$. With $d=\delta_1=\delta_2=0$, Daemi and Scaduto's complex computing $\widehat I(K;\Q)$ is
$C_{*-1}\oplus\Q[x]$, with $x$ acting by $v$ on the first summand~\cite[\S4.3, Lemma~4.11]{DSEquiv}, so
$(C,v)$ is the $\Q[x]$-torsion of $\widehat I(K;\Q)$, up to a shift of degree, and it is an invariant of
$K$~\cite[Thm.~1.4]{DSEquiv}.
\end{proof}

\begin{proof}[Proof of Corollary~\ref{cor:graded}]
Substitute $R_1=\lceil R/2\rceil$ and $R_3=\lfloor R/2\rfloor$ into $\beta$.
\end{proof}

\begin{example}\label{ex:45}
For $T(4,5)$, $\sigma=-8$ and $N_1=N_3=2$, so $\tC$ has rank vector $(3,2,2,2)$ in degrees $0,1,2,3$, and
since $\sigma\equiv0\pmod4$ no shift is needed to compare with $\Inat$. Kronheimer and Mrowka showed that
$\Inat(T(4,5))$ has rank vector $(2,1,2,2)$~\cite[\S11]{KMfilt}; see also~\cite[\S12.5]{HHK1}. A
differential of degree $-1$ with these chain and homology ranks has a single nonzero component, of rank
one, from degree $1$ to degree $0$. The maps of degree $1\to0$ are $\delta_1$ and $v_{13}$, and $\delta_1=0$
because $h(T(4,5))=0$. Hence $v_{13}$ has rank one and $v_{31}=0$. This is the instanton counterpart of
the single pillowcase bigon of $T(4,5)$, which also lowers the grading from $1$ to $0$ and avoids the
reducible generator~\cite[\S11.8]{HHK2}.
\end{example}

\begin{remark}\label{rem:vnotinv}
The rank of $v$ itself is not an invariant when $\delta_1\neq0$: for $T(3,4)$ and $T(3,5)$, the
representatives with $v=0$ and with $v=w\,\delta_1$ ($w\in C_3$, $w\ne0$) are S-isomorphic, by
Theorem~\ref{thm:split}(1) with $P=1$. Daemi and Scaduto observe this freedom in the choice of
splitting~\cite[Rem.~4.4]{DSEquiv}.
\end{remark}

\begin{remark}\label{rem:3n}
For $n\equiv4,5\pmod6$, $h(T(3,n))=1$. Write $\Upsilon_n=\Upsilon_{T(3,n)}(1)$ and
$\sigma_n=\sigma(T(3,n))$. By~\cite[Prop.~2.2]{FK}, $\Upsilon_{n+3}=\Upsilon_n+\Upsilon_4$ for $n\ge1$, and
$\Upsilon_4=-2$, since $h(T(3,4))=1$~\cite[\S4.1, Prop.~10]{DScaduto} and $\sigma_4=-6$. Since also
$\sigma_{n+6}=\sigma_n-8$~\cite[Thm.~5.2]{GLM}, $h(T(3,n))=\Upsilon_n-\sigma_n/2$ depends only on $n$ modulo
$6$. Finally $\Upsilon_5=\Upsilon_2+\Upsilon_4=-3$ and $\sigma_5=-8$, so $h(T(3,5))=1$. By Theorem~\ref{thm:split}(3), $v_{31}=0$ and $v_{13}$ vanishes on $\ker\delta_1$, so
$v_{13}=w\,\delta_1$ for some $w\in C_3$, and the S-isomorphism $(1,-w)$ makes $v=0$: up to S-isomorphism
over $\Q$ the differential of $T(3,n)$ is $\delta_1$ alone.
\end{remark}

\section{Comparison with the pillowcase}\label{sec:pillow}

Hedden, Herald and Kirk decompose $T(p,q)$ along a Conway sphere bounding a ball in which the knot is a
trivial tangle. The trivial tangle carries the earring and gives a figure-eight curve $L_0$ in the
pillowcase; the other tangle gives, after a small holonomy perturbation, an immersed curve
$L_1$~\cite[\S\S10--11]{HHK2}. Table~\ref{tab:pillow} compares their computations with
Theorem~\ref{thm:split}. After the grading is lowered by $\sigma$, a bigon into $r_+$ corresponds to
$\delta_1$ (degree $1\to0$); any other bigon from grading $\sigma+3$ to $\sigma+2$ corresponds to $v_{31}$
(degree $3\to2$); and any other bigon from grading $\sigma+1$ to $\sigma$ corresponds to
$v_{13}|_{\ker\delta_1}$ (degree $1\to0$).

\begin{table}[ht]
\centering\small
\begin{tabular}{c|c|c|c|l|l}
knot & $R$ & $h$ & $\sigma\bmod4$ & bigons in~\cite{HHK2}, gradings of~\cite{HHK2} & component of $\td$\\ \hline
$T(3,4)$ & $1$ & $1$ & $2$ & $x_1^-\to r_+$, $3\to2$ (\S11.6) & $\delta_1$\\
$T(3,5)$ & $1$ & $1$ & $0$ & $x_1^-\to r_+$, $1\to0$ (\S11.7) & $\delta_1$\\
$T(4,5)$ & $1$ & $0$ & $0$ & one bigon avoiding $r_+$, $1\to0$ (\S11.8) & $v_{13}$\\
$T(4,7)$ & $2$ & $1$ & $2$ & $x_3^-\to r_+$, $3\to2$; $x_1^-\to x_2^+$, $1\to0$ (\S11.9) & $\delta_1$; $v_{31}$\\
$T(5,11)$ & $4$ & $0$ & $0$ & four bigons on closed components (\S11.1) & $v_{31}$, $v_{13}$, rank $2$ each\\
\end{tabular}
\medskip
\caption{Torus knots with nonzero pillowcase differential in~\cite{HHK2}. The gradings in the fifth column
are those of~\cite{HHK2}; subtracting $\sigma$ gives the degrees of the components in the last column,
which are determined by Theorem~\ref{thm:split} and Example~\ref{ex:45}. In each case a bigon ends at
$r_+$ exactly when $h>0$, and the rank of the pillowcase differential equals $R$. The knots $T(3,7)$,
$T(5,7)$, $T(5,12)$ and $T(5,17)$ of~\cite[\S11]{HHK2} have $R=0$ and no differential.}
\label{tab:pillow}
\end{table}

For $T(5,11)$, the curve $L_1$ consists of an arc, which contributes $r_+$, and four closed
components~\cite[\S11.1]{HHK2}. The four bigons lie on two of the closed components, two on each, and
cancel $8=25-17$ of the $25$ generators in pairs; as a linear map the differential has rank $4=R$. On each
of these components the two bigons run from some $x_i^-$ to some $x_j^+$ with $j\neq i$~\cite[\S3.9]{HHK2},
as Conjecture~\ref{conj:pillow}(a) requires; the absolute gradings of closed components are fixed by
convention in~\cite{HHK2}, so this row does not test the split between $v_{31}$ and $v_{13}$.
For $T(4,7)$, Theorem~\ref{thm:split} gives $\operatorname{rank}v_{31}=1$ and $v_{13}|_{\ker\delta_1}=0$,
which is what the two bigons show. (Hedden, Herald and Kirk name the second pair $x_2^+,x_1^-$; its
direction $x_1^-\to x_2^+$ follows from the gradings.) Over $\mathbb F_2$, the bigon into $r_+$ for $T(4,7)$
also predicts that $\delta_1$ is nonzero modulo $2$, which is not known. The pillowcase complex depends on
the perturbation: for larger perturbation the pair joined by the bigon into $r_+$ cancels and the complex
of $T(3,4)$ becomes perfect~\cite[\S11.10, Prop.~11.1]{HHK2}. This is consistent with $h(T(3,4))=1$: once
that pair has cancelled, the complex no longer carries the S-structure of Conjecture~\ref{conj:pillow},
whereas $h$ is an invariant of the S-complex. The conjecture therefore concerns the regime of small
perturbations. With the earring on the other tangle the pillowcase homology of $T(4,5)$ has the wrong
rank~\cite[\S10.3]{CHKK}, so the choice of decomposition matters as well, and so does the choice of $(r,s)$
within the family of~\cite[\S10]{HHK2}~\cite{Wpill}. A proof would amount to a
chain-level form of the knot Atiyah--Floer correspondence of~\cite{HHK2,CHKK} for torus knots;
Cazassus, Herald, Kirk and Kotelskiy already suggest viewing the earring as a counterpart of the
equivariant variable $x$~\cite[\S11.6]{CHKK}.

\begin{thebibliography}{HHK18}

\bibitem[Bat14]{Batson} J.~Batson, \emph{Nonorientable slice genus can be arbitrarily large}, Math.\ Res.\
Lett.\ \textbf{21} (2014), 423--436.

\bibitem[BO10]{BO} M.~Borodzik, K.~Oleszkiewicz, \emph{On the signatures of torus knots}, Bull.\ Pol.\
Acad.\ Sci.\ Math.\ \textbf{58} (2010), 167--177.

\bibitem[CHKK22]{CHKK} G.~Cazassus, C.~Herald, P.~Kirk, A.~Kotelskiy, \emph{The correspondence induced on
the pillowcase by the earring tangle}, J.\ Topol.\ \textbf{15} (2022), 2472--2543.

\bibitem[DS24a]{DScaduto} A.~Daemi, C.~Scaduto, \emph{Chern--Simons functional, singular instantons, and
the four-dimensional clasp number}, J.\ Eur.\ Math.\ Soc.\ \textbf{26} (2024), 2127--2190.

\bibitem[DS24b]{DSEquiv} A.~Daemi, C.~Scaduto, \emph{Equivariant aspects of singular instanton Floer
homology}, Geom.\ Topol.\ \textbf{28} (2024), 4057--4190.

\bibitem[FK17]{FK} P.~Feller, D.~Krcatovich, \emph{On cobordisms between knots, braid index, and the
Upsilon-invariant}, Math.\ Ann.\ \textbf{369} (2017), 301--329.

\bibitem[GLM81]{GLM} C.~McA.~Gordon, R.~A.~Litherland, K.~Murasugi, \emph{Signatures of covering links},
Canad.\ J.\ Math.\ \textbf{33} (1981), 381--394.

\bibitem[HHK14]{HHK1} M.~Hedden, C.~Herald, P.~Kirk, \emph{The pillowcase and perturbations of traceless
representations of knot groups}, Geom.\ Topol.\ \textbf{18} (2014), 211--287.

\bibitem[HHK18]{HHK2} M.~Hedden, C.~Herald, P.~Kirk, \emph{The pillowcase and traceless representations of
knot groups II: a Lagrangian--Floer theory in the pillowcase}, J.\ Symplectic Geom.\ \textbf{16} (2018),
721--815.

\bibitem[JVC21]{JVC} S.~Jabuka, C.~A.~Van Cott, \emph{On a nonorientable analogue of the Milnor
conjecture}, Algebr.\ Geom.\ Topol.\ \textbf{21} (2021), 2571--2625.

\bibitem[KM10a]{KMAlex} P.~B.~Kronheimer, T.~S.~Mrowka, \emph{Instanton Floer homology and the
Alexander polynomial}, Algebr.\ Geom.\ Topol.\ \textbf{10} (2010), 1715--1738.

\bibitem[KM10b]{KMsut} P.~B.~Kronheimer, T.~S.~Mrowka, \emph{Knots, sutures, and excision}, J.\
Differential Geom.\ \textbf{84} (2010), 301--364.

\bibitem[KM11a]{KM1} P.~B.~Kronheimer, T.~S.~Mrowka, \emph{Knot homology groups from instantons}, J.\
Topol.\ \textbf{4} (2011), 835--918.

\bibitem[KM11b]{KM2} P.~B.~Kronheimer, T.~S.~Mrowka, \emph{Khovanov homology is an unknot-detector},
Publ.\ Math.\ IH\'ES \textbf{113} (2011), 97--208.

\bibitem[KM14]{KMfilt} P.~B.~Kronheimer, T.~S.~Mrowka, \emph{Filtrations on instanton homology}, Quantum
Topol.\ \textbf{5} (2014), 61--97.

\bibitem[Lim10]{Lim} Y.~Lim, \emph{Instanton homology and the Alexander polynomial}, Proc.\ Amer.\
Math.\ Soc.\ \textbf{138} (2010), 3759--3768.

\bibitem[Lit79]{Litherland} R.~A.~Litherland, \emph{Signatures of iterated torus knots}, in
\emph{Topology of Low-Dimensional Manifolds}, Lecture Notes in Math.\ \textbf{722}, Springer, 1979,
71--84.

\bibitem[LY22]{LiYe} Z.~Li, F.~Ye, \emph{Instanton Floer homology, sutures, and Heegaard diagrams}, J.\
Topol.\ \textbf{15} (2022), 39--107.

\bibitem[Mor14]{Moree} P.~Moree, \emph{Numerical semigroups, cyclotomic polynomials, and Bernoulli numbers},
Amer.\ Math.\ Monthly \textbf{121} (2014), 890--902.

\bibitem[OSS17]{OSS} P.~S.~Ozsv\'ath, A.~I.~Stipsicz, Z.~Szab\'o, \emph{Concordance homomorphisms from knot
Floer homology}, Adv.\ Math.\ \textbf{315} (2017), 366--426.


\bibitem[Wue26a]{WuebbenI} B.~J.~Wuebben, \emph{Traceless $\mathrm{SU}(2)$ characters and $\mathbb{Z}/4$ instanton
gradings for two-bridge and $(3,n)$-torus knots}, arXiv:2607.26095; current version at \url{https://github.com/bwuebben/pillowcase-bounding-cochain}.

\bibitem[Wue26b]{WuebbenII} B.~J.~Wuebben, \emph{Instanton and pillowcase homology of the $(-2,3,q)$ pretzel
knots}, arXiv:2607.26096 (earlier versions under a different title); current version at \url{https://github.com/bwuebben/pillowcase-bounding-cochain}.

\bibitem[Wue26c]{Wpill} B.~J.~Wuebben, \emph{Pillowcase Floer homology of the torus knots $T(3,n)$}, preprint,
2026, \url{https://github.com/bwuebben/pillowcase-bounding-cochain}.

\end{thebibliography}

\bigskip
\noindent{\sc Bernd Johannes Wuebben, New York, NY,} \texttt{wuebben@gmail.com}

\end{document}
